\section{AMC 12 12A 2011}
        \begin{enumerate}
        	\item (\textit{AMC 12 12A 2011}) A cell phone plan costs $\$20$ dollars each month, plus $5$ cents per text message sent, plus $10$ cents for each minute used over $30$ hours. In January Michelle sent $100$ text messages and talked for $30.5$ hours. How much did she have to pay?
$\textbf{(A)}\ 24.00 \qquad \textbf{(B)}\ 24.50 \qquad \textbf{(C)}\ 25.50 \qquad \textbf{(D)}\ 28.00 \qquad \textbf{(E)}\ 30.00$

 \textbf{Answer: }D



\textbf{Solution}

The base price of Michelle's cell phone plan is $20$ dollars. If she sent $100$ text messages and it costs $5$ cents per text, then she must have spent $500$ cents for texting, or $5$ dollars. She talked for $30.5$ hours, but $30.5-30$ will give us the amount of time exceeded past 30 hours.$30.5-30=.5$ hours $=30$ minutes.Since the price for phone calls is $10$ cents per minute, the additional amount Michelle has to pay for phone calls is $30*10=300$ cents, or $3$ dollars.Adding $20+5+3$ dollars $\boxed{\$28 \textbf{ (D)}}$.

 


	\item (\textit{AMC 12 12A 2011}) There are $5$ coins placed flat on a table according to the figure. What is the order of the coins from top to bottom?

                    \begin{figure}[H]
                        \centering
                        \includegraphics[width=0.5\textwidth]{Problem/AMC_12_12A_2011_2_Problem_Figure_1.png}
                    \end{figure}
                    
$\textbf{(A)}\ (C, A, E, D, B) \qquad \textbf{(B)}\ (C, A, D, E, B) \qquad \textbf{(C)}\ (C, D, E, A, B) \qquad \textbf{(D)}\ (C, E, A, D, B) \qquad \\ \textbf{(E)}\ (C, E, D, A, B)$

 \textbf{Answer: }E



\textbf{Solution}

By careful inspection and common sense, the answer is $\textbf{(E)}$.

 \textbf{Note} Note that this image is quite similar to a topological map. This means that the coin that appears closest to us is at the top and the coin that appears farthest from us is at the bottom.

 


	\item (\textit{AMC 12 12A 2011}) A small bottle of shampoo can hold $35$ milliliters of shampoo, whereas a large bottle can hold $500$ milliliters of shampoo. Jasmine wants to buy the minimum number of small bottles necessary to completely fill a large bottle. How many bottles must she buy?


 $\textbf{(A)}\ 11 \qquad \textbf{(B)}\ 12 \qquad \textbf{(C)}\ 13 \qquad \textbf{(D)}\ 14 \qquad \textbf{(E)}\ 15$

 \textbf{Answer: }E



\textbf{Solution 1}

To find how many small bottles we need, we can simply divide $500$ by $35$. This simplifies to $\frac{100}{7}=14 \frac{2}{7}.$ Since the answer must be an integer greater than $14$, we have to round up to $15$ bottles, or $\boxed{\textbf{E}}$




\textbf{Solution 2}

We double $35$ to get $70.$ We see that $70\cdot7=490,$ which is very close to $500.$ Thus, $2\cdot7+1=\boxed{\text{(E)} 15}$ bottles are enough. ~Technodoggo

 


	\item (\textit{AMC 12 12A 2011}) At an elementary school, the students in third grade, fourth grade, and fifth grade run an average of $12$, $15$, and $10$ minutes per day, respectively. There are twice as many third graders as fourth graders, and twice as many fourth graders as fifth graders. What is the average number of minutes run per day by these students?


 $\textbf{(A)}\ 12 \qquad \textbf{(B)}\ \frac{37}{3} \qquad \textbf{(C)}\ \frac{88}{7} \qquad \textbf{(D)}\ 13 \qquad \textbf{(E)}\ 14$

 \textbf{Answer: }C



\textbf{Solution}

Let us say that there are $f$ fifth graders. According to the given information, there must be $2f$ fourth graders and $4f$ third graders. The average time run by each student is equal to the total amount of time run divided by the number of students. This gives us $\frac{12\cdot 4f + 15\cdot 2f + 10\cdot f}{4f + 2f + f} = \frac{88f}{7f} = \frac{88}{7} \Rightarrow \boxed{C}$

 


 If you want to simplify the problem even more, just imagine/assume that only $1$ fifth grader existed. Then you can simply get rid of the variables.

 


	\item (\textit{AMC 12 12A 2011}) Last summer $30\%$ of the birds living on Town Lake were geese, $25\%$ were swans, $10\%$ were herons, and $35\%$ were ducks. What percent of the birds that were not swans were geese?


 $\textbf{(A)}\ 20 \qquad \textbf{(B)}\ 30 \qquad \textbf{(C)}\ 40 \qquad \textbf{(D)}\ 50 \qquad \textbf{(E)}\ 60$

 \textbf{Answer: }C



\textbf{Solution}

To simplify the problem, WLOG, let us say that there were a total of $100$ birds. The number of birds that are not swans is $75$. The number of geese is $30$. Therefore the percentage is just $\frac{30}{75} \times 100 = 40 \Rightarrow \boxed{C}$



	\item (\textit{AMC 12 12A 2011}) The players on a basketball team made some three-point shots, some two-point shots, and some one-point free throws. They scored as many points with two-point shots as with three-point shots. Their number of successful free throws was one more than their number of successful two-point shots. The team's total score was $61$ points. How many free throws did they make?


 $\textbf{(A)}\ 13 \qquad \textbf{(B)}\ 14 \qquad \textbf{(C)}\ 15 \qquad \textbf{(D)}\ 16 \qquad \textbf{(E)}\ 17$

 \textbf{Answer: }A



\textbf{Solution 1}

For the points made from two-point shots and from three-point shots to be equal, the numbers of made shots are in a $3:2$ ratio. Therefore, assume they made $3x$ and $2x$ two- and three- point shots, respectively, and thus $3x+1$ free throws. The total number of points is 
 \[2 \times (3x) + 3 \times (2x) + 1 \times (3x+1) = 15x+1\]

 Set that equal to $61$, we get $x = 4$, and therefore the number of free throws they made $3 \times 4 + 1 = 13 \Rightarrow \boxed{A}$




\textbf{Solution 2}

Let $x$ be the number of free throws. Then the number of points scored by two-pointers is $2(x-1)$ and the same goes for three-pointers because they scored the same number of points with twos and threes. Thus, our equation is $x+4(x-1) = 61 \Rightarrow x=13$, giving us $\boxed{(A)}$ for an answer.

 



\textbf{Solution 3}

We let $a$ be the number of $2$-point shots, $b$ be the number of $3$-point shots, and $x$ be the number of free throws. We are looking for $x.$ We know that $2a=3b$, and that $x=a+1$. Also, $2a+3b+1x=61$. We can see

 
 \begin{align*} a&=x-1 \\ 2a &= 2x-2 \\ 3a &= 2x-2. \\ \end{align*}


 Plugging this into $2a+3b+1x=61$, we see

 
 \begin{align*} 2x-2+2x-2+x &= 61 \\ 5x-4 &= 61 \\ 5x &= 65 \\ x &= \boxed{\textbf{(A) }13}. \end{align*}




	\item (\textit{AMC 12 12A 2011}) A majority of the $30$ students in Ms. Demeanor's class bought pencils at the school bookstore. Each of these students bought the same number of pencils, and this number was greater than $1$. The cost of a pencil in cents was greater than the number of pencils each student bought, and the total cost of all the pencils was $17.71$. What was the cost of a pencil in cents?


 $\textbf{(A)}\ 7 \qquad \textbf{(B)}\ 11 \qquad \textbf{(C)}\ 17 \qquad \textbf{(D)}\ 23 \qquad \textbf{(E)}\ 77$

 \textbf{Answer: }B



\textbf{Solution 1}

The total cost of the pencils can be found by $(\text{students}\cdot\text{pencils purchased by each}\cdot\text{price of each pencil})$.

 Since $1771$ is the product of three sets of values, we can begin with prime factorization, since it gives some insight into the values: $7, 11, 23$. Since $(C)$ has none of these factors, it can be eliminated immediately, leaving $(A)$, $(B)$, $(D)$, and $(E)$.

 Beginning with $(A) 7$, we see that the number of pencils purchased by each student must be either $11$ or $23$. However, the problem states that the price of each pencil must exceed the number of pencils purchased, so we can eliminate this.

 Continuing with $(B) 11$, we can conclude that the only case that fulfills the restrictions is that there are $23$ students who each purchased $7$ such pencils, so the answer is $\boxed{B}$. We can apply the same logic to $(E)$ as we applied to $(A)$ if one wants to make doubly sure.

 



\textbf{Solution 2}

We know the total cost of the pencils can be found by $(\text{students}\cdot\text{pencils purchased by each}\cdot\text{price of each pencil})$.

 Using prime factorization like in the solution above, we see that there is only one combination of three whole numbers whose product is equal to $1771$: $7, 11, 23$ (without using $1$). So we know that $7, 11$, and $23$ must be the number of students, the number of pencils purchased by each student and the price of each pencil in cents.

 We know that $23$ must be the number of students, as it is the only number that makes up the majority of 30.

 We pick the greater of the remaining numbers for the price of each pencil in cents, which is $11$.

 Therefore, our answer is $\boxed{B}$.

 -JinhoK

 



\textbf{Solution 3 (essentially Solutions 1 and 2 but cleaned with variables)}

We let $s$ be the number of students that bought pencils at the bookstore, $c$ be the cost of each pencil in cents, and $n$ the number of pencils each student bought. Thus, we are looking for $c.$ Since a majority of the students in the class bought pencils at the bookstore, $s>\dfrac{30}2=15.$ (s>15) We also know that $n>1$ and $c>n>1.$ Finally, $s\cdot c\cdot n=1771.$ We can factor $1771$ as $7\cdot11\cdot23.$ Since $n>15$ and $c>n>1,$ none of them can be $1,$ and therefore $c,n,$ and $s$ are $7,11,$ and $23$ in some order. We know that $n>15,$ so $n$ must be $23.$ $c>n,$ so $c=11$ and $n=7.$ Thus, $c=\boxed{\text{(B)}~11}$ cents.~Technodoggo

 


	\item (\textit{AMC 12 12A 2011}) In the eight term sequence $A$, $B$, $C$, $D$, $E$, $F$, $G$, $H$, the value of $C$ is $5$ and the sum of any three consecutive terms is $30$. What is $A+H$?


 $\textbf{(A)}\ 17 \qquad \textbf{(B)}\ 18 \qquad \textbf{(C)}\ 25 \qquad \textbf{(D)}\ 26 \qquad \textbf{(E)}\ 43$

 \textbf{Answer: }C



\textbf{Solution 1}

Let $A=x$. Then from $A+B+C=30$, we find that $B=25-x$. From $B+C+D=30$, we then get that $D=x$. Continuing this pattern, we find $E=25-x$, $F=5$, $G=x$, and finally $H=25-x$. So $A+H=x+25-x=25 \rightarrow \boxed{\textbf{C}}$




\textbf{Solution 2}

Given that the sum of 3 consecutive terms is 30, we have$(A+B+C)+(C+D+E)+(F+G+H)=90$ and $(B+C+D)+(E+F+G)=60$

 It follows that $A+B+C+D+E+F+G+H=85$ because $C=5$.

 Subtracting, we have that $A+H=25\rightarrow \boxed{\textbf{C}}$.

 


 



\textbf{Solution 3 (the tedious one)}

From the given information, we can deduce the following equations:

 $A+B=25, B+D=25, D+E=25, D+E+F=30, E+F+G =30$, and $F+G+H=30$.

 We can then cleverly manipulate the equations above by adding and subtracting them to be left with the answer.

 $(A+B)-(B+D)=25-25 \implies (A-D)=0$
$(A-D)+(D+E)=0+25 \implies (A+E)=25$
$(A+E)-(E+F+G)=25-30 \implies (A-F-G)=-5$ (Notice how we don't use $D+E+F=30$)

 $(A-F-G)+(F+G+H)=-5+30 \implies (A+H)=25$

 Therefore, we have $A+H=25 \rightarrow \boxed{\textbf{C}}$




\textbf{Solution 4 (the cheap one)}

Since all of the answer choices are constants, it shouldn't matter what we pick $A$ and $B$ to be, so let $A = 20$ and $B = 5$. Then $D = 30 - B -C = 20$, $E = 30 - D - C = 5$, $F = 30 - D - E =5$, and so on until we get $H = 5$. Thus $A + H = \boxed{\mathbf{(C)}25}$




\textbf{Solution 5}

Notice that the period of the sequence is $3$ as given. (If this isn't clear we can show an example: $A+B+C=B+C=D$ $\Leftrightarrow$ $A=D$). Then $A=D$ and $H=B$, so $A+H=D+B=D+B+C-C=30-5=\boxed{\textbf{(C)}\;25}$.

 



\textbf{Solution 6}

Since the period of the sequence is $3$ 
 \[A=D=G, B=E=H, C=F\]$A+B+C=30$ because any three consecutive terms sum to $30$.Since $C=5$ 
 \[A+B=25\]Since $B=H$ 
 \[A+B=A+H=\fbox{(C) 25}\]

 \textbf{Note} Something useful to shorten a lot of the solutions above is to notice 
 \[5 + D + E = D + E + F\] so F = 5

 


	\item (\textit{AMC 12 12A 2011}) At a twins and triplets convention, there were $9$ sets of twins and $6$ sets of triplets, all from different families. Each twin shook hands with all the twins except his/her siblings and with half the triplets. Each triplet shook hands with all the triplets except his/her siblings and with half the twins. How many handshakes took place?


 $\textbf{(A)}\ 324 \qquad \textbf{(B)}\ 441 \qquad \textbf{(C)}\ 630 \qquad \textbf{(D)}\ 648 \qquad \textbf{(E)}\ 882$

 \textbf{Answer: }B



\textbf{Solution}

There are $18$ total twins and $18$ total triplets. Each of the twins shakes hands with the $16$ twins not in their family and $9$ of the triplets, a total of $25$ people. Each of the triplets shakes hands with the $15$ triplets not in their family and $9$ of the twins, for a total of $24$ people. Dividing by two to accommodate the fact that each handshake was counted twice, we get a total of $\frac{1}{2}\times 18 \times (25+24) = 9 \times 49 = 441 \rightarrow \boxed{\textbf{B}}$



	\item (\textit{AMC 12 12A 2011}) A pair of standard $6$-sided dice is rolled once. The sum of the numbers rolled determines the diameter of a circle. What is the probability that the numerical value of the area of the circle is less than the numerical value of the circle's circumference?


 $\textbf{(A)}\ \frac{1}{36} \qquad \textbf{(B)}\ \frac{1}{12} \qquad \textbf{(C)}\ \frac{1}{6} \qquad \textbf{(D)}\ \frac{1}{4} \qquad \textbf{(E)}\ \frac{5}{18}$

 \textbf{Answer: }B



\textbf{Solution}

For the circumference to be greater than the area, we must have $\pi d > \pi \left( \frac{d}{2} \right) ^2$, or $d<4$. Now since $d$ is determined by a sum of two dice, the only possibilities for $d$ are thus $2$ and $3$. In order for two dice to sum to $2$, they must both show a value of $1$. The probability of this happening is $\frac{1}{6} \times \frac{1}{6} = \frac{1}{36}$. In order for two dice to sum to $3$, one must show a $1$ and the other must show a $2$. Since this can happen in two ways, the probability of this event occurring is $2 \times \frac{1}{6} \times \frac{1}{6} = \frac{2}{36}$. The sum of these two probabilities now gives the final answer: $\frac{1}{36} + \frac{2}{36} = \frac{3}{36} = \frac{1}{12} \rightarrow \boxed{\textbf{B}}$



	\item (\textit{AMC 12 12A 2011}) Circles $A, B,$ and $C$ each have radius 1. Circles $A$ and $B$ share one point of tangency. Circle $C$ has a point of tangency with the midpoint of $\overline{AB}.$ What is the area inside circle $C$ but outside circle $A$ and circle $B?$
$\textbf{(A)}\ 3 - \frac{\pi}{2} \qquad \textbf{(B)}\ \frac{\pi}{2} \qquad \textbf{(C)}\  2 \qquad \textbf{(D)}\ \frac{3\pi}{4} \qquad \textbf{(E)}\ 1+\frac{\pi}{2}$

 \textbf{Answer: }C



\textbf{Solution 1}


                    \begin{figure}[H]
                        \centering
                        \includegraphics[width=0.5\textwidth]{Solution/AMC_12_12A_2011_11_Solution_Figure_1.png}
                    \end{figure}
                    

 The requested area is the area of $C$ minus the area shared between circles $A$, $B$ and $C$.

 Let $M$ be the midpoint of $\overline{AB}$ and $D$ be the other intersection of circles $C$ and $B$.

 The area shared between $C$, $A$ and $B$ is $4$ of the regions between arc $\widehat {MD}$ and line $\overline{MD}$, which is (considering the arc on circle $B$) a quarter of the circle $B$ minus $\triangle MDB$:

 $\frac{\pi r^2}{4}-\frac{bh}{2}$
$b = h = r = 1$

 (We can assume this because $\angle DBM$ is 90 degrees, since $CDBM$ is a square, due to the application of the tangent chord theorem at point $M$)

 So the area of the small region is

 $\frac{\pi}{4}-\frac{1}{2}$

 The requested area is area of circle $C$ minus 4 of this area:

 $\pi 1^2 - 4\left(\frac{\pi}{4}-\frac{1}{2}\right) = \pi - \pi + 2 = 2$
$\boxed{\textbf{C}}$.

 



\textbf{Solution 2}


                    \begin{figure}[H]
                        \centering
                        \includegraphics[width=0.5\textwidth]{Solution/AMC_12_12A_2011_11_Solution_Figure_1.png}
                    \end{figure}
                    

 We can move the area above the part of the circle above the segment $EF$ down, and similarly for the other side. Then, we have a square, whose diagonal is $2$, so the area is then just $\left(\frac{2}{\sqrt{2}}\right)^2 = \boxed{\textbf{2 = C}}$.

 



\textbf{Solution 3}

We can approach the problem by analyzing the overlapping regions of the figures. Let the regions associated with the triangles be denoted as follows: let $A$ and $B$ represent the two small congruent triangles, and let $C$ represent the primary region from which we subtract the overlapping components to find our desired area.

 We can find the two small congruent overlapping segments from triangles $A$ and $B$ and subtract them from the area of $C$ to find our answer. We can see that half of each overlapping segment is a quarter circle minus an isosceles right triangle of area $\frac{1}{2}$. This means that the area of half of the intersecting region is:
 \[\frac{\pi}{4} - \frac{1}{2}\]

 Now we need to multiply this by $2$ to get the area of the whole intersecting region:
 \[2 \left( \frac{\pi}{4} - \frac{1}{2} \right) = \frac{\pi}{2} - 1\]

 We now multiply this by $2$ to find the total area of the $2$ intersecting regions, which gives:
 \[2 \left( \frac{\pi}{2} - 1 \right) = \pi - 2\]

 Subtracting this total intersecting area from the area of region $C$ (which is $\pi$), we find the area of the region of $C$ that we are looking for:
 \[\pi - (\pi - 2) = 2\]

 Thus, the final area is $\boxed{2}$, which corresponds to option (D).

 


	\item (\textit{AMC 12 12A 2011}) A power boat and a raft both left dock $A$ on a river and headed downstream. The raft drifted at the speed of the river current. The power boat maintained a constant speed with respect to the river. The power boat reached dock $B$ downriver, then immediately turned and traveled back upriver. It eventually met the raft on the river 9 hours after leaving dock $A.$ How many hours did it take the power boat to go from $A$ to $B?$
$\textbf{(A)}\ 3 \qquad \textbf{(B)}\ 3.5 \qquad \textbf{(C)}\  4 \qquad \textbf{(D)}\ 4.5 \qquad \textbf{(E)}\ 5$

 \textbf{Answer: }D



\textbf{Solution 1}

WLOG, let the speed of the river be 0. This is allowed because the problem never states that the speed of the current has to have a magnitude greater than 0. In this case, when the powerboat travels from $A$ to $B$, the raft remains at $A$. Thus the trip from $A$ to $B$ takes the same time as the trip from $B$ to the raft. Since these times are equal and sum to $9$ hours, the trip from $A$ to $B$ must take half this time, or $4.5$ hours. The answer is thus $\boxed{\textbf{D}}$.

 Remark: This is equivalent to viewing the problem from the raft's perspective.

 



\textbf{Solution 2}

What's important in this problem is to consider everything in terms of the power boat and the raft, since that is how the problem is given to us. Think of the blue arrow as the power boat and the red arrow as the raft in the following three diagrams, which represent different time intervals of the problem.

 
                    \begin{figure}[H]
                        \centering
                        \includegraphics[width=0.5\textwidth]{Solution/AMC_12_12A_2011_12_Solution_Figure_1.png}
                    \end{figure}
                    

 Thinking about the distance covered as their distances with respect to each other, they are $0$ distance apart in the first diagram when they haven't started to move yet, some distance $d$ apart in the second diagram when the power boat reaches $B$, and again $0$ distance apart in the third diagram when they meet. Therefore, with respect to each other, the boat and the raft cover a distance of $d$ on the way there, and again cover a distance of $d$ on when drawing closer. This makes sense, because from the 1st diagram to the second, the raft moves in the same direction as the boat, while from the 2nd to the 3rd, the boat and raft move in opposite directions.

 Let $b$ denote the speed of the power boat (only the power boat, not factoring in current) and $r$ denote the speed of the raft, which, as given by the problem, is also equal to the speed of the current. Thus, from $A$ to $B$, the boat travels at a velocity of $b+r$, and on the way back, travels at a velocity of $-(b-r)=r-b$, since the current aids the boat on the way there, and goes against the boat on the way back. With respect to the raft then, the boat's velocity from $A$ to $B$ becomes $(r+b)-r=b$, and on the way back it becomes $(r-b)-r=-b$. Since the boat's velocities with respect to the raft are exact opposites, $b$ and $-b$, we therefore know that the boat and raft travel apart from each other at the same rate that they travel toward each other.

 From this, we have that the boat travels a distance $d$ at rate $b$ with respect to the raft both on the way to $B$ and on the way back. Thus, using $\dfrac{distance}{speed}=time$, we have $\dfrac{2d}{b}=9\text{ hours}$, and to see how long it took to travel half the distance, we have $\dfrac{d}{b}=4.5\text{ hours}\implies\boxed{\textbf{D}}$




\textbf{Solution 3}

Let $t$ be the time it takes the power boat to go from $A$ to $B$ in hours, $r$ be the speed of the river current (and thus also the raft), and $p$ to be the speed of the power boat with respect to the river. 

 Using $d = rt$, the raft covers a distance of $9r$, the distance from $A$ to $B$ is $(p + r)t$, and the distance from $B$ to where the raft and power boat met up is $(9 - t)(p - r)$. 

 Then, $9r + (9 - t)(p - r) = (p + r)t$. Solving for $t$, we get $t = 4.5$, which is $\boxed{\textbf{D}}$.

 


 


	\item (\textit{AMC 12 12A 2011}) Triangle $ABC$ has side-lengths $AB = 12, BC = 24,$ and $AC = 18.$ The line through the incenter of $\triangle ABC$ parallel to $\overline{BC}$ intersects $\overline{AB}$ at $M$ and $\overline{AC}$ at $N.$ What is the perimeter of $\triangle AMN?$
$\textbf{(A)}\ 27 \qquad \textbf{(B)}\ 30 \qquad \textbf{(C)}\  33 \qquad \textbf{(D)}\ 36 \qquad \textbf{(E)}\ 42$

 \textbf{Answer: }B



\textbf{Solution 1}

Let $O$ be the incenter of $\triangle{ABC}$. Because $\overline{MO} \parallel \overline{BC}$ and $\overline{BO}$ is the angle bisector of $\angle{ABC}$, we have 

 
 \[\angle{MBO} = \angle{CBO} = \angle{MOB} = \frac{1}{2}\angle{MBC}\]

 It then follows due to alternate interior angles and base angles of isosceles triangles that $MO = MB$. Similarly, $NO = NC$. The perimeter of $\triangle{AMN}$ then becomes
 \begin{align*} AM + MN + NA &= AM + MO + NO + NA \\ &= AM + MB + NC + NA \\ &= AB + AC \\ &= 30 \rightarrow \boxed{(B)} \end{align*}





\textbf{Solution 2}

Let $O$ be the incenter. $AO$ is the angle bisector of $\angle MAN$. Let the angle bisector of $\angle BAC$ meets $BC$ at $P$ and the angle bisector of $\angle ABC$ meets $AC$ at $Q$. By applying both angle bisector theorem and Menelaus' theorem,

 


 $\frac{AO}{OP} \times \frac{BP}{BC} \times \frac{CQ}{QA} = 1$
$\frac{AO}{OP} \times \frac{12}{30} \times \frac{24}{12} = 1$
$\frac{AO}{OP}=\frac{5}{4}$
$\frac{AO}{AP}=\frac{5}{9}$

 
Perimeter of $\triangle AMN = \frac{12+24+18}{9} \times 5 = 30  \rightarrow \boxed{(B)}$




\textbf{Solution 3}

Like in other solutions, let $O$ be the incenter of $\triangle ABC$. Let $AO$ intersect $BC$ at $D$. By the angle bisector theorem, $\frac{BD}{DC} = \frac{AB}{AC} = \frac{12}{18} = \frac{2}{3}$. Since $BD+DC = 24$, we have $\frac{BD}{24-BD} = \frac{2}{3}$, so $3BD = 48 - 2BD$, so $BD = \frac{48}{5} = 9.6$. By the angle bisector theorem on $\triangle ABD$, we have $\frac{DO}{OA} = \frac{BD}{BA} = \frac{4}{5} = 0.8$, so $\frac{DA}{OA} = 1 + \frac{DO}{OA} = \frac{9}{5} = 1.8$, so $\frac{AO}{AD} = \frac{5}{9}$. Because $\triangle AMN \sim \triangle ABC$, the perimeter of $\triangle AMN$ must be $\frac{5}{9} (12 + 18 + 24) = 30$, so our answer is $\boxed{\textbf{(B)}\ 30}$.

 Another way to find $\frac{AO}{OD}$ is to use mass points. Assign a mass of 24 to $A$, a mass of 18 to $B$, and a mass of 12 to $C$. Then $D$ has mass 30, so $\frac{AO}{OD} = \frac{30}{24} = \frac{5}{4} = 1.25$.

 



\textbf{Solution 4}

We know that the ratio of the perimeter of $\triangle AMN$ and $\triangle ABC$ is the ratio of their heights, and finding the two heights is pretty easy. Note that the height from $A$ to $BC$ is $\frac{9\sqrt{15}}{4}$ from Herons and then $A=\frac{bh}{2}$ and also that the height from $A$ to $MN$ is simply the height from $A$ to $BC$ minus the inradius. We know the area and the semiperimeter so $r=\frac{A}{s}$ which gives us $r=\sqrt{15}$. Now we know that the altitude from $A$ to $MN$ is $\frac{5\sqrt{15}}{4}$ so the ratios of the heights from $A$ for $\triangle AMN$ and $\triangle ABC$ is $\frac{5}{9}$. Thus the perimeter of $\triangle AMN$ is $\frac{5}{9} \times 54 = 30$ so our answer is $\boxed{B}$

 -srisainandan6

 


	\item (\textit{AMC 12 12A 2011}) Suppose $a$ and $b$ are single-digit positive integers chosen independently and at random. What is the probability that the point $(a,b)$ lies above the parabola $y=ax^2-bx$?


 $\textbf{(A)}\ \frac{11}{81} \qquad \textbf{(B)}\ \frac{13}{81} \qquad \textbf{(C)}\ \frac{5}{27} \qquad \textbf{(D)}\ \frac{17}{81} \qquad \textbf{(E)}\ \frac{19}{81}$

 \textbf{Answer: }E



\textbf{Solution}

If $(a,b)$ lies above the parabola, then $b$ must be greater than $y(a)$. We thus get the inequality $b>a^3-ba$. Solving this for $b$ gives us $b>\frac{a^3}{a+1}$. Now note that $\frac{a^3}{a+1}$ constantly increases when $a$ is positive. Then since this expression is greater than $9$ when $a=4$, we can deduce that $a$ must be less than $4$ in order for the inequality to hold, since otherwise $b$ would be greater than $9$ and not a single-digit integer. The only possibilities for $a$ are thus $1$, $2$, and $3$.

 For $a=1$, we get $b>\frac{1}{2}$ for our inequality, and thus $b$ can be any integer from $1$ to $9$.

 For $a=2$, we get $b>\frac{8}{3}$ for our inequality, and thus $b$ can be any integer from $3$ to $9$.

 For $a=3$, we get $b>\frac{27}{4}$ for our inequality, and thus $b$ can be any integer from $7$ to $9$.

 Finally, if we total up all the possibilities we see there are $19$ points that satisfy the condition, out of $9 \times 9 = 81$ total points. The probability of picking a point that lies above the parabola is thus $\frac{19}{81} \rightarrow \boxed{\textbf{E}}$



	\item (\textit{AMC 12 12A 2011}) The circular base of a hemisphere of radius $2$ rests on the base of a square pyramid of height $6$. The hemisphere is tangent to the other four faces of the pyramid. What is the edge-length of the base of the pyramid?


 $\textbf{(A)}\ 3\sqrt{2} \qquad \textbf{(B)}\ \frac{13}{3} \qquad \textbf{(C)}\ 4\sqrt{2} \qquad \textbf{(D)}\ 6 \qquad \textbf{(E)}\ \frac{13}{2}$

 \textbf{Answer: }A



\textbf{Solution 1}

Let $ABCDE$ be the pyramid with $ABCD$ as the square base. Let $O$ and $M$ be the center of square $ABCD$ and the midpoint of side $AB$ respectively. Lastly, let the hemisphere be tangent to the triangular face $ABE$ at $P$.

 Notice that $\triangle EOM$ has a right angle at $O$. Since the hemisphere is tangent to the triangular face $ABE$ at $P$, $\angle EPO$ is also $90^{\circ}$. Hence $\triangle EOM$ is similar to $\triangle EPO$.

 $\frac{OM}{2} = \frac{6}{EP}$
$OM = \frac{6}{EP} \times 2$
$OM = \frac{6}{\sqrt{6^2 - 2^2}} \times 2 = \frac{3\sqrt{2}}{2}$

 The length of the square base is thus $2 \times \frac{3\sqrt{2}}{2} = 3\sqrt{2} \rightarrow \boxed{\textbf{A}}$




\textbf{Solution 2 (Answer Choices)}

Consider a cross section of the pyramid such that it is a hemisphere inscribed inside of a triangle. Let $A$ be the vertex furthest away from the hemisphere. Let $P$ be a point where the hemisphere is tangent. Let $O$ be the centre of the hemisphere. Triangle $AOP$ is a right triangle, since $P$ is ninety degrees. Apply the pythagorean formula to find side $AP$ as $4\sqrt{2}.$ Since we know the answer must have $\sqrt{2}$ somewhere, we can remove choices B, D and E from selection. Choice C is simply the length of $AP$, which cannot be the desired length. Thus, the answer must be $\boxed{\textbf{A}}$.

 



\textbf{Solution 3(Cross Sections)}

The cross section of this entire figure is just an isosceles triangle with height $6$ and a semicircle tangent to the congruent sides of the triangle and its base is on the base of the triangle. This semicircle has radius $2$. Now draw the line connecting the center of this semicircle and the point of tangency with one of the sides. This will have length $2$. Let the hypotenuse of this new right triangle we just created have length $x$. Therefore, the altitude of the isosceles triangle with height $6$ will be similar to this triangle with hypotenuse $x$. Hence, the hypotenuse of this longer triangle is just $\sqrt{36 + x^2}$. By similar triangles

 $\frac{2}{6} = \frac{x}{\sqrt{36 + x^2}} \implies \sqrt{36 + x^2} = 3x \implies 36 + x^2 = 9x^2 \implies x^2 = \frac{36}{8} = \frac{9}{2} \implies x = \frac{3\sqrt{2}}{2}$. Since this is only a half of the base of the isosceles triangle, we double $x$ to get

 $\boxed{\textbf{A}}$.

 


	\item (\textit{AMC 12 12A 2011}) Each vertex of convex polygon $ABCDE$ is to be assigned a color. There are $6$ colors to choose from, and the ends of each diagonal must have different colors. How many different colorings are possible?


 $\textbf{(A)}\ 2520 \qquad \textbf{(B)}\ 2880 \qquad \textbf{(C)}\ 3120 \qquad \textbf{(D)}\ 3250 \qquad \textbf{(E)}\ 3750$

 \textbf{Answer: }C



\textbf{Solution}

We can do some casework when working our way around the pentagon from $A$ to $E$. At each stage, there will be a makeshift diagram.

 1.) For $A$, we can choose any of the 6 colors.

 \begin{verbatim}
		A : 6
\end{verbatim}

2.) For $B$, we can either have the same color as $A$, or any of the other 5 colors. We do this because each vertex of the pentagon is affected by the 2 opposite vertices, and $D$ will be affected by both $A$ and $B$.

 \begin{verbatim}
	A : 6 B:1        B:5
\end{verbatim}

3.) For $C$, we cannot have the same color as $A$. Also, we can have the same color as $B$ ($E$ will be affected), or any of the other 4 colors. Because $C$ can't be the same as $A$, it can't be the same as $B$ if $B$ is the same as $A$, so it can be any of the 5 other colors.

 \begin{verbatim}
	A : 6 B:1        B:5 C:5     C:4   C:1
\end{verbatim}

4.) $D$ is affected by $A$ and $B$. If they are the same, then $D$ can be any of the other 5 colors. If they are different, then $D$ can be any of the (6-2)=4 colors.

 \begin{verbatim}
	A : 6 B:1        B:5 C:5     C:4   C:1 D:5     D:4   D:4
\end{verbatim}

5.) $E$ is affected by $B$ and $C$. If they are the same, then $E$ can be any of the other 5 colors. If they are different, then $E$ can be any of the (6-2)=4 colors.

 \begin{verbatim}
	A : 6 B:1        B:5 C:5     C:4   C:1 D:5     D:4   D:4 E:4     E:4   E:5
\end{verbatim}

6.) Now, we can multiply these three paths and add them:$(6\times1\times5\times5\times4)+(6\times5\times4\times4\times4)+(6\times5\times1\times4\times5) =600+1920+600=3120$

 7.) Our answer is $C$!

 



\textbf{Solution 2}

Right off the bat, we can analyze three things: 

 
1.) There can only be two of the same color on the pentagon.

 2.) Any pair of the same color can only be next to each other on the pentagon.

 3.) There can only be two different pairs of same colors on the pentagon at once.

 
Now that we know this, we can 

 


	\item (\textit{AMC 12 12A 2011}) Circles with radii $1$, $2$, and $3$ are mutually externally tangent. What is the area of the triangle determined by the points of tangency?


 $\textbf{(A)}\ \frac{3}{5} \qquad \textbf{(B)}\ \frac{4}{5} \qquad \textbf{(C)}\ 1 \qquad \textbf{(D)}\ \frac{6}{5} \qquad \textbf{(E)}\ \frac{4}{3}$

 \textbf{Answer: }D



\textbf{Solution 1}


                    \begin{figure}[H]
                        \centering
                        \includegraphics[width=0.5\textwidth]{Solution/AMC_12_12A_2011_17_Solution_Figure_1.png}
                    \end{figure}
                    

 The centers of these circles form a 3-4-5 triangle, which has an area equal to 6.

 The areas of the three triangles determined by the center and the two points of tangency of each circle are, using Triangle Area by Sine,

 $\frac{1}{2} \cdot 1 \cdot 1 \cdot 1 = \frac{1}{2}$
$\frac{1}{2} \cdot 2 \cdot 2 \cdot \frac{4}{5} = \frac{8}{5}$
$\frac{1}{2} \cdot 3 \cdot 3 \cdot \frac{3}{5} = \frac{27}{10}$

 which add up to $4.8$. The area we're looking for is the large 3-4-5 triangle minus the three smaller triangles, or $6 - 4.8 = 1.2 = \frac{6}{5} \rightarrow \boxed{(D)}$.

 



\textbf{Solution 2 (Analytical)}

Let $O_1,O_2,O_3$ be the centers of the circles with radii $1,2,3$. Notice that the points of tangency of the $3$ circles are also the points of tangency of the incircle of $\triangle O_1O_2O_3$. Using the radius of an Incircle formula, $r = \frac{A}{S}$ where $S$ is the semi-perimeter, and noting that $\triangle O_1O_2O_3$ is a 3-4-5 right triangle, we see that, 
 \[r = \frac{3 \cdot 4/2}{\frac{3+4+5}{2}} = 1.\]Now we set $\triangle O_1O_2O_3$ on the coordinate plane with $O_1=(0,0),O_2=(3,0), O_3 = (0,4)$. So the incenter lies on $(1,1)$. Let the points of tangency of $\triangle O_1O_2O_3$ with it's incenter are $A,B,C$ with $A$ on $O_1O_2$, $B$ on $O_1O_3$, and $C$ on $O_2O_3$. We have that $A = (1,0), B=(0,1)$. Since the line defined by $C$ and the incenter is perpendicular to $O_2O_3$ who has equation $y = \frac{-4}{3}x + 4$, we have it's equation as $(y-1) = \frac{3}{4}(x-1) \rightarrow y = \frac{3}{4}x + \frac{1}{4}$. We have the intersection of the $2$ lines at, 
 \begin{align*}\frac{-4}{3}x + 4 &= \frac{3}{4}x + \frac{1}{4} \\ \frac{15}{4} &= \frac{25}{12}x \\ x &= \frac{9}{5} \\ y &= \frac{-4}{3} \frac{9}{5} + 4 = \frac{8}{5}.\end{align*}
Here we can use Shoelace Theorem on the points $(\frac{9}{5}, \frac{8}{5}),(0,1),(1,0)$ we get our areas as $\frac{6}{5} \rightarrow \boxed{(D)}.$



	\item (\textit{AMC 12 12A 2011}) Suppose that $\left|x+y\right|+\left|x-y\right|=2$. What is the maximum possible value of $x^2-6x+y^2$?


 $\textbf{(A)}\ 5 \qquad \textbf{(B)}\ 6 \qquad \textbf{(C)}\ 7 \qquad \textbf{(D)}\ 8 \qquad \textbf{(E)}\ 9$

 \textbf{Answer: }D



\textbf{Solution 1}

Plugging in some values, we see that the  graph of the equation $|x+y|+|x-y| = 2$ is a square bounded by $x= \pm 1$ and $y = \pm 1$.

 Notice that $x^2 - 6x + y^2 = (x-3)^2 + y^2 - 9$ means the square of the distance from a point $(x,y)$ to point $(3,0)$ minus 9. To maximize that value, we need to choose the point in the feasible region farthest from point $(3,0)$, which is $(-1, \pm 1)$. Either one, when substituting into the function, yields $\boxed{\textbf{(D)}\ 8}$.

 



\textbf{Solution 2}

Since the equation $|x+y|+|x-y| = 2$ is dealing with absolute values, the following could be deduced: $(x+y)+(x-y)=2$,$(x+y)-(x-y)=2$, $-(x+y)+(x-y)=2$, and $-(x+y)-(x-y)=2$. Simplifying would give $x=1$, $y=1$, $y=-1$, and $x=-1$. In $x^2-6x+y^2$, we care most about $-6x,$ since both $x^2$ and $y^2$ are non-negative. To maximize $-6x$, though, $x$ would have to be -1. Therefore, when $x=-1$ and $y=-1$ or $y=1$, the equation evaluates to $\boxed{\textbf{(D)}\ 8}$.

 


	\item (\textit{AMC 12 12A 2011}) At a competition with $N$ players, the number of players given elite status is equal to $2^{1+\lfloor \log_{2} (N-1) \rfloor}-N$. Suppose that $19$ players are given elite status. What is the sum of the two smallest possible values of $N$?


 $\textbf{(A)}\ 38 \qquad \textbf{(B)}\ 90 \qquad \textbf{(C)}\ 154 \qquad \textbf{(D)}\ 406 \qquad \textbf{(E)}\ 1024$

 \textbf{Answer: }C



\textbf{Solution 1}

We start with $2^{1+\lfloor\log_{2}(N-1)\rfloor}-N = 19$. After rearranging, we get $\lfloor\log_{2}(N-1)\rfloor = \log_{2} \left(\frac{N+19}{2}\right)$.

 Since $\lfloor\log_{2}(N-1)\rfloor$ is a positive integer, $\frac{N+19}{2}$ must be in the form of $2^{m}$ for some positive integer $m$. From this fact, we get $N=2^{m+1}-19$.

 If we now check integer values of N that satisfy this condition, starting from $N=19$, we quickly see that the first values that work for $N$ are $45$ and $109$, that is, $2^6-19$ and $2^7 -19$, giving values of $5$ and $6$ for $m$, respectively. Adding up these two values for $N$, we get $45 + 109 = 154 \rightarrow \boxed{\textbf{C}}$




\textbf{Solution 2}

We examine the value that $2^{1+\lfloor\log_{2}(N-1)\rfloor}$ takes over various intervals. The $\lfloor\log_{2}(N-1)\rfloor$ means it changes on each multiple of 2, like so:

 2 --> 1

 3 - 4 --> 2

 5 - 8 --> 3

 9 - 16 --> 4

 From this, we see that $2^{1+\lfloor\log_{2}(N-1)\rfloor} - N$ is the difference between the next power of 2 above $2^{\lfloor\log_{2}(N-1)\rfloor}$ and $N$. We are looking for $N$ such that this difference is 19. The first two $N$ that satisfy this are $45 = 64-19$ and $109=128-19$ for a final answer of $45 + 109 = 154 \rightarrow \boxed{\textbf{C}}$




\textbf{Solution 3 (using the answer choices)}

Note that each $N$ is $19$ less than a power of $2$. So, the answer will be $38$ less than the sum of $2$ powers of $2$. Adding $38$ to each answer, we get $76$, $128$, $192$, $444$, and $1062$. Obviously we can take out $76$ and $1062$. Also, $128$ will not work because two powers of two will never sum to another power of $2$ (unless they are equal, which is a contradiction to the question). So, we have $192$ and $444$. Note that $444 = 1 + 443 = 2 + 442 = 4 + 440 = 8 + 436 = 16 + 428 = 32 + 412$, etc. We quickly see that $444$ will not work, leaving $192$ which corresponds to $\boxed{\textbf{C}}$. We can also confirm that this works because $192 = 128 + 64 = 2^7 + 2^6$.

 



\textbf{Solution 4 (removing the log)}

In order to fix the exponent and get rid of the logarithm term, let $N = 2^m + k + 1$, with $0 \leq k < 2^m$. Doing so, we see that $\lfloor \log_2{N - 1} \rfloor = m$, which turns our given relation into
 \[2^m = 20 + k,\]for which the solutions of the form $(m, k)$, $(5, 12)$ and $(6, 44)$, follow trivially. Adding up the two values of $N$ gives us $32 + 12 + 1 + 64 + 44 + 1 = 154$, so the answer is $\boxed{\textbf{C}}$.

 


	\item (\textit{AMC 12 12A 2011}) Let $f(x)=ax^2+bx+c$, where $a$, $b$, and $c$ are integers. Suppose that $f(1)=0$, $50<f(7)<60$, $70<f(8)<80$, $5000k<f(100)<5000(k+1)$ for some integer $k$. What is $k$?


 $\textbf{(A)}\ 1 \qquad \textbf{(B)}\ 2 \qquad \textbf{(C)}\ 3 \qquad \textbf{(D)}\ 4 \qquad \textbf{(E)}\ 5$

 \textbf{Answer: }C



\textbf{Solution 1}

From $f(1) = 0$, we know that $a+b+c = 0$.

 From the first inequality, we get $50 < 49a+7b+c < 60$. Subtracting $a+b+c = 0$ from this gives us $50 < 48a+6b < 60$, and thus $\frac{25}{3} < 8a+b < 10$. Since $8a+b$ must be an integer, it follows that $8a+b = 9$.

 Similarly, from the second inequality, we get $70 < 64a+8b+c < 80$. Again subtracting $a+b+c = 0$ from this gives us $70 < 63a+7b < 80$, or $10 < 9a+b < \frac{80}{7}$. It follows from this that $9a+b = 11$.

 We now have a system of three equations: $a+b+c = 0$, $8a+b = 9$, and $9a+b = 11$. Solving gives us 

 


	\item (\textit{AMC 12 12A 2011}) Let $f_{1}(x)=\sqrt{1-x}$, and for integers $n \geq 2$, let $f_{n}(x)=f_{n-1}(\sqrt{n^2 - x})$. If $N$ is the largest value of $n$ for which the domain of $f_{n}$ is nonempty, the domain of $f_{N}$ is $\{ c\}$. What is $N+c$?


 $\textbf{(A)}\ -226 \qquad \textbf{(B)}\ -144 \qquad \textbf{(C)}\ -20 \qquad \textbf{(D)}\ 20 \qquad \textbf{(E)}\ 144$

 \textbf{Answer: }A



\textbf{Solution 1}

The domain of $f_{1}(x)=\sqrt{1-x}$ is defined when $x\leq1$. 
 \[f_{2}(x)=f_{1}\left(\sqrt{4-x}\right)=\sqrt{1-\sqrt{4-x}}\] 

 Applying the domain of $f_{1}(x)$ and the fact that square roots must be positive, we get $0\leq\sqrt{4-x}\leq1$. Simplifying, the domain of $f_{2}(x)$ becomes $3\leq x\leq4$. 

 Repeat this process for $f_{3}(x)=\sqrt{1-\sqrt{4-\sqrt{9-x}}}$ to get a domain of $-7\leq x\leq0$.  

 For $f_{4}(x)$, since square roots must be nonnegative, we can see that the negative values of the previous domain will not work, so $\sqrt{16-x}=0$. Thus we now arrive at $16$ being the only number in the of domain of $f_4 x$ that defines $x$. However, since we are looking for the largest value for $n$ for which the domain of $f_{n}$ is nonempty, we must continue checking until we arrive at a domain that is empty. 

 We continue with $f_{5}(x)$ to get a domain of $\sqrt{25-x}=16 \implies x=-231$. Since square roots cannot be negative, this is the last nonempty domain. We add to get $5-231=\boxed{\textbf{(A)}\ -226}$. 

 



\textbf{Solution 2}

We start with smaller values. Notice that $f_4(x) = \sqrt{1-\sqrt{4-\sqrt{9-\sqrt{16-x}}}}$. Notice that the mess after $\sqrt{1-(mess)}$ must be greater than 0, since it's a square root, and less than 1, since otherwise the inside of the larger square root on the outside would be negative. 

 
Continuing, we get that $\sqrt{16-x} = 0$, which means $x=16$ is the only value in the domain of $f_4(x)$. Now we move on to $f_5(x)$. The only change with $f_5(x)$ is replacing the $x$ from $f_4(x)$ with $\sqrt{25-x}$. Since we had $x=16$ in $f_4(x)$, in $f_5(x)$, $\sqrt{25-x} = 16$, forcing $x=-231$.

 
Clearly, we can't move on from here, since $f_6(x)$ would replace $x$ with $\sqrt{36-x}$, and we would need $-231 = \sqrt{36-x}$, but a square root can never be negative, so $N=5$, $c=-231$, and the answer is $5-231 = \boxed{\textbf{(A) }-226}$.

 -skibbysiggy

 


	\item (\textit{AMC 12 12A 2011}) Let $R$ be a unit square region and $n \geq 4$ an integer. A point $X$ in the interior of $R$ is called \textit{n-ray partitional} if there are $n$ rays emanating from $X$ that divide $R$ into $n$ triangles of equal area. How many points are $100$-ray partitional but not $60$-ray partitional?


 $\textbf{(A)}\ 1500 \qquad \textbf{(B)}\ 1560 \qquad \textbf{(C)}\ 2320 \qquad \textbf{(D)}\ 2480 \qquad \textbf{(E)}\ 2500$

 \textbf{Answer: }C



\textbf{Solution 1}

There must be four rays emanating from $X$ that intersect the four corners of the square region. Depending on the location of $X$, the number of rays distributed among these four triangular sectors will vary. We start by finding the corner-most point that is $100$-ray partitional (let this point be the bottom-left-most point). 

 We first draw the four rays that intersect the vertices. At this point, the triangular sectors with bases as the sides of the square that the point is closest to both do not have rays dividing their areas. Therefore, their heights are equivalent since their areas are equal. The remaining $96$ rays are divided among the other two triangular sectors, each sector with $48$ rays, thus dividing these two sectors into $49$ triangles of equal areas. 

 Let the distance from this corner point to the closest side be $a$ and the side of the square be $s$. From this, we get the equation $\frac{a\times s}{2}=\frac{(s-a)\times s}{2}\times\frac1{49}$. Solve for $a$ to get $a=\frac s{50}$. Therefore, point $X$ is $\frac1{50}$ of the side length away from the two sides it is closest to. By moving $X$ $\frac s{50}$ to the right, we also move one ray from the right sector to the left sector, which determines another $100$-ray partitional point. We can continue moving $X$ right and up to derive the set of points that are $100$-ray partitional. 

 In the end, we get a square grid of points each $\frac s{50}$ apart from one another. Since this grid ranges from a distance of $\frac s{50}$ from one side to $\frac{49s}{50}$ from the same side, we have a $49\times49$ grid, a total of $2401$ $100$-ray partitional points. To find the overlap from the $60$-ray partitional, we must find the distance from the corner-most $60$-ray partitional point to the sides closest to it. Since the $100$-ray partitional points form a $49\times49$ grid, each point $\frac s{50}$ apart from each other, we can deduce that the $60$-ray partitional points form a $29\times29$ grid, each point $\frac s{30}$ apart from each other. To find the overlap points, we must find the common divisors of $30$ and $50$ which are $1, 2, 5,$ and $10$. Therefore, the overlapping points will form grids with points $s$, $\frac s{2}$, $\frac s{5}$, and $\frac s{10}$ away from each other respectively. Since the grid with points $\frac s{10}$ away from each other includes the other points, we can disregard the other grids. The total overlapping set of points is a $9\times9$ grid, which has $81$ points. Subtract $81$ from $2401$ to get $2401-81=\boxed{\textbf{(C)}\ 2320}$.

 



\textbf{Solution 2}

Position the square region $R$ so that the bottom-left corner of the square is at the origin. Then define $s$ to be the sidelength of $R$ and $X$ to be the point $(rs, qs)$, where $0<r,q<1$. 

 
There must be four rays emanating from $X$ that intersect the four corners of $R$. The areas of the four triangles formed by these rays are then $A_1=\frac{qs\times s}{2}$, $A_2=\frac{(s-rs)\times s}{2}$, $A_3=\frac{(s-qs)\times s}{2}$, and $A_4=\frac{rs\times s}{2}$. 

 
If a point is $n$-ray partitional, then there exist positive integers $a, b, c, d$ such that $a+b+c+d=n$ and $\frac{A_1}{a}=\frac{A_2}{b}=\frac{A_3}{c}=\frac{A_4}{d}$. Substituting in our formulas for $A_1$, $A_2$, $A_3$, and $A_4$ and canceling equal terms, we get 
 \[\frac{q}{a}=\frac{1-r}{b}=\frac{1-q}{c}=\frac{r}{d}.\] 

 
Taking $\frac{q}{a}=\frac{1-q}{c}$ and solving for $q$, we get $q=\frac{a}{a+c}$, and taking $\frac{1-r}{b}=\frac{r}{d}$ and solving for $r$, we get $r=\frac{d}{b+d}$. Finally, from $\frac{q}{a}=\frac{r}{d}$, we have $qd=ar$ $\Leftrightarrow$ $\frac{ad}{a+c}=\frac{ad}{b+d}$ $\Leftrightarrow$ $a+c=b+d$. 

 
So for a point $X$ to be $100$-ray partitional, $a+b+c+d=100$, so $a+c=b+d=50$. $X$ must then be of the form $\left(\frac{d}{50}s, \frac{a}{50}s\right)$. Since $X$ is in the interior of $R$, $a$ and $d$ can be any positive integer from $1$  to $49$ (with $b$ and $c$ just equaling $50-d$ and $50-a$, respectively). Thus, there are $49\times 49=2401$ points that are $100$-ray partitional.

 
However, the problem asks for points that are not only $100$-ray partitional but also not $60$-ray partitional. Points that are $60$-ray partitional are of the form $\left(\frac{m}{30}s, \frac{n}{30}s\right)$, where $m$ and $n$ are also positive integers. We count the number of points $\left(\frac{d}{50}s, \frac{a}{50}s\right)$ that can also be written in this form. For a given $d$, $\frac{d}{50}=\frac{m}{30}$ if and only if $m=\frac{3}{5}d$, and likewise with $a$ and $n$. We can then see that a point is both $100$-ray partitional and $60$-ray partitional if and only if $a$ and $d$ are both divisible by $5$. There are $9$ integers between $1$ and $49$ that are divisible by $5$, so out of our $2401$ points that are $100$-ray partitional, $9\times 9=81$ are also $60$-ray partitional.

 
Our answer then is just $2401-81=\boxed{\textbf{(C)}\ 2320}$.

 



\textbf{Solution 3}

For the sake of simplicity, let $R$ be a $60 \times 60$ square and set the bottom-left point as the origin. Then, $R$ has vertices: 
 \[(0,0), (60,0), (60,60), (0,60).\]

 Now, let a point in the square have coordinates $(x, y).$

 In order for the point to be $100-$ ray partitional, we must be able to make $100$ triangles with area $60^2/100 = 36.$ For it to not be $60$-ray partitional, we cannot make $60$ triangles with area $60^2/60 = 60.$

 When we draw such a triangle, it's base will always be on one of the sides of the square. If it is on the bottom side, then the height must be $y$. So, the base of each triangle must be $\frac{72}{y}.$ So, there will be $\frac{60}{\frac{72}{y}} = \frac{60y}{72}$ total triangles on that side.

 If the triangle is on the right side of the square, then the height has to be $60 - x.$ So, the base will be $\frac{72}{60-x}$ and there will be $\frac{60 (60-x)}{72}$ triangles.

 If the triangle is on the top, the height will be $60-y$, the base will be $\frac{72}{60-y}$ and there will be $\frac{60 (60-y)}{72}$ triangles along this side.

 The triangles on the left will have height $x$, base $\frac{72}{x}$ and $\frac{60x}{72}$ such triangles must exist. 

 Simplifying, we get $\frac{5y}{6}, \frac{5 (60-x)}{6}, \frac{5(60-y)}{6}, \frac{5x}{6}$ triangles on the respective sides of the square. 

 All of these numbers must be integers. Let $x = 60a$ and $y = 60b$ where $0 < a, b < 1.$

 
Then, the amounts become: 
 \[50a, 50b,  50 - 50a, 50 - 50b.\]

 As long as $50a$ and $50b$ are integers, $50 - 50a$ and $50 - 50b$ will also be integers. 

 For this to happen, $a$ and $b$ must be of the form $\frac{x}{50}$ where $1 \le x \le 49.$ So, we have a total of $49^2 = 2401$ points that are $100$-ray partitional. 

 Now, we must calculate the number of $100$- ray partitional points that are also $60$-ray partitional. 

 Using a method similar to before, if a point is $60$-ray partitional, then we must be able to make $30a, 30b, 30 - 30a, 30 -30b$ triangles on the different sides. 

 So, $30a$ and $30b$ must be integers. This means $a$ and $b$ have to be of the form $\frac{y}{30}$ where $0<y<30.$

 If a point is both $100$-ray partitional and $60$-ray partitional, then it can be written as 
 \[\frac{x}{50} =  \frac{y}{30}\]. Note that whenever $x$ is divisible by $5$, a $y$ will always exist to satisfy the above equation.

 So, $x$ has to be in the range $(0, 50)$ and must be divisible by $5$. There are $9$ possibilities, namely $5, 10, 15, 20, 25, 30, 35, 40, 45.$ 

 The $x$-coordinate of the point and the $y$-coordinate of the point can each use any of these $9$ possibilities, giving $9^2 = 81$ numbers that are both $100$-ray partitional and $60$-ray partitional.

 Overall, we are left with $2401 - 81 = \boxed{2320}$ solutions.

 



\textbf{Solution 4 (not scary)}

If $X$ is $100$-ray partitional, and there are $n$ triangles on a side, then the total area of those $n$ triangles will be $\frac{n}{100}$. Since the side has length $1$, then the length of the altitude from $X$ to that side will be $\frac{n}{50}$.

 
There will be 1 unique way to split this side into bases of triangles with height $\frac{n}{50}$, and also 1 unique split into triangles on the opposite side. Since $n$ can range from $[1,49]$, there will be $49$ different distances from $1$ side, and $49$ different distances from an adjacent side, for a total of $49\cdot 49 = 2401$ $100-ray$ partitional points. 

 
Similarly, for a point to be $60$-ray partitional, the distance from the sides should be of the form $\frac{k}{30}$. For this point to also be $100$-ray partitional, it must also be of the form $\frac{z}{50}$. Clearly, the only possible values of $k$ are $(3,6,9,12,15,18,21,24,27)$ for a total of 9 ways, so there are $9\cdot 9 = 81$ points that are $60$-ray partitional and $100$-ray partitional.

 
Subtracting $81$ from $2401$, we get $\boxed{\textbf{(C) } 2320}$. 

 -skibbysiggy

 


	\item (\textit{AMC 12 12A 2011}) Let $f(z)= \frac{z+a}{z+b}$ and $g(z)=f(f(z))$, where $a$ and $b$ are complex numbers. Suppose that $\left| a \right| = 1$ and $g(g(z))=z$ for all $z$ for which $g(g(z))$ is defined. What is the difference between the largest and smallest possible values of $\left| b \right|$?


 $\textbf{(A)}\ 0 \qquad \textbf{(B)}\ \sqrt{2}-1 \qquad \textbf{(C)}\ \sqrt{3}-1 \qquad \textbf{(D)}\ 1 \qquad \textbf{(E)}\ 2$

 \textbf{Answer: }C



\textbf{Solution 1 (plug and chug)}

By algebraic manipulations, we obtain 
 \[h(z)=g(g(z))=f(f(f(f(z))))=\frac{Pz+Q}{Rz+S}\] where 
 \[P=(a+1)^2+a(b+1)^2\]
 \[Q=a(b+1)(b^2+2a+1)\]
 \[R=(b+1)(b^2+2a+1)\]
 \[S=a(b+1)^2+(a+b^2)^2\]In order for $h(z)=z$, we must have $R=0$, $Q=0$, and $P=S$. 

 $R=0$ implies $b=-1$ or $b^2+2a+1=0$.

 $Q=0$ implies $a=0$, $b=-1$, or $b^2+2a+1=0$.

 $P=S$ implies $b=\pm1$ or $b^2+2a+1=0$. 

 Since $|a|=1\neq 0$, in order to satisfy all 3 conditions we must have either $b=\pm1$ or $b^2+2a+1=0$. In the first case $|b|=1$. 

 For the latter case note that 
 \[|b^2+1|=|-2a|=2\]
 \[2=|b^2+1|\leq |b^2|+1\]and hence, 
 \[1\leq|b|^2\Rightarrow1\leq |b|\]. On the other hand, 
 \[2=|b^2+1|\geq|b^2|-1\]so, 
 \[|b^2|\leq 3\Rightarrow0\leq |b|\leq \sqrt{3}\]. Thus $1\leq |b|\leq \sqrt{3}$. Hence the maximum value for $|b|$ is $\sqrt{3}$ while the minimum is $1$ (which can be achieved for instance when $|a|=1,|b|=\sqrt{3}$ or $|a|=1,|b|=1$ respectively). Therefore the answer is $\boxed{\textbf{(C)}\ \sqrt{3}-1}$.

 \textbf{Shortcut} We only need $Q$ in $f^4(z)=g^2(z)=\frac{Pz+\textcolor{red}{Q}}{Rz+S}$.

 Set $Q=0$: $a(b+1)\left(b^2+2a+1\right)=0$. Since $|a|=1$, either $b+1=0$ or $b^2+2a+1=0$.

 $b+1=0\rightarrow b=-1$ so $|b|=1$.

 $b^2+2a+1=0\rightarrow b^2=-1-2a$. This is a circle in the complex plane centered at $(-1,0)$ with radius $2$ since $|a|=1$. The maximum distance from the origin is $3$ at $(-3,0)$ and similarly the minimum distance is $1$ at $(1,0)$. So $1\le |b^2|\le 3\rightarrow 1\le |b|\le \sqrt{3}$.

 Both solutions give the same lower bound, $1$. So the range is $\sqrt{3}-1=\boxed{\textbf{(C) }\sqrt{3}-1}$.

 



\textbf{Solution 2}

note: the "begin{align*}" environment is weird on aops, so if the two things that use it below are duplicated in front of each other, just reload the page or something.

 $\textbf{Lemma.}$ Let $h(x)=\frac{ax+b}{cx+d}$ and $h\left (h(x)\right )=x$ for all such $x$ where $h(h(x))$ is defined. Then either $d=-a$ or $h(x)=x$.

 This lemma is one that is typically known by people who could make it this far into this AMC 12 anyway (USAMO qualifiers likely). 

 $\textit{Proof.}$ We plug in $h\left (h(x)\right )$, and equating it to $x$, we see\begin{align*}h\left (h(x)\right ) &=h\left (\frac{ax+b}{cx+d}\right ) \\ &= \frac{a\cdot \frac{ax+b}{cx+d}+b}{c\cdot \frac{ax+b}{cx+d}+d} \\ &= \frac{a(ax+b)+b(cx+d)}{c(ax+b)+d(cx+d)} \\ &= \frac{(a^2+bc)x+(ab+bd)}{(ac+cd)x+(bc+d^2)}=x.\end{align*}
Then if we clear the denominator, it rearranges into 
 \[(a^2+bc)x+(ab+bd)=(ac+cd)x^2+(bc+d^2)x\implies (ac+cd)x^2+(d^2-a^2)x-(ab+bd)=0.\] Since this equation is true for all $x$ (since now the fractions are gone), this implies that the coefficients of this quadratic are all equal to $0$. Focus on the condition $d^2-a^2=0$.

 $\textbf{Case 1:}$ $d+a=0$

 Then all the other coefficients automatically become zero, so $d+a=0\implies d=-a$ is a sufficient condition.

 $\textbf{Case 2:}$ $d-a=0$ (and $d+a\neq 0$)

 Then in $ab+bd=0$, we divide both sides by $a+d$ to see that $b=0$. Similarly, in $ac+cd=0$, we divide to also see that $c=0$. Then since $d=a$, 
 \[h(x)=\frac{ax+b}{cx+d}=\frac{ax+0}{0x+a}=x.\] Therefore, in this case, it is required that $h(x)=x$, which now works. This proves our desired claim. $\qquad \qquad \square$

 Note that we are asking when $g(g(z))=z$. So either $g(z)=z$, or in the expansion of $g(z)=f(f(z))$, the two corresponding coefficients are negations of each other. 

 $\textbf{Case 1:}$ $g(z)=z$

 Then $f(f(z))=z$, so either $f(z)=z$ or the two coefficients are negative of each other. Clearly $f(z)\neq z$, since the numerator and denominator are both degree $1$ polynomials. So we would require the coefficients, namely $1$ and $b$, to be negatives of each other. Then $b=-1$ is a possibility here, giving $|b|=1$.

 $\textbf{Case 2:}$ The two coefficients are negative of each other.

 It can be found that \begin{align*}f(f(x))&=f\left (\frac{z+a}{z+b}\right ) \\ &=\frac{\frac{z+a}{z+b}+a}{\frac{z+a}{z+b}+b} \\ &=\frac{z+a+a(z+b)}{z+a+b(z+b)} \\ &= \frac{(a+1)z+(ab+a)}{(b+1)z+(b^2+a)}.\end{align*}
 By our lemma, we would require $a+1=-(b^2+a)$, which rearranges into $b^2=-1-2a$. But we are given that $|a|=1$. The triangle inequality for complex numbers states that $|a|+|b|\geq |a+b|$, where equality happens when the arguments of $a$ and $b$ are equal. Here, we can see that 
 \[|-1-2a|\leq |-1|+|-2a|=1+2=3,\] but we can also find that 
 \[|-1-2a|+1=|-1-2a|+|1|\geq |-2a|=2\implies |-1-2a|\geq 1.\] Therefore, because $|b^2|=|-1-2a|$, then $1\leq |b|^2\leq 3$, giving us $1\leq |b|\leq \sqrt{3}$.

 When we combine the two cases, we find that the possible range of $|b|$ is $|b|\in [1,\sqrt{3}]$, so the answer is $\boxed{\textbf{(C) }\sqrt{3}-1}$.

 


	\item (\textit{AMC 12 12A 2011}) Consider all quadrilaterals $ABCD$ such that $AB=14$, $BC=9$, $CD=7$, and $DA=12$. What is the radius of the largest possible circle that fits inside or on the boundary of such a quadrilateral?


 $\textbf{(A)}\ \sqrt{15} \qquad \textbf{(B)}\ \sqrt{21} \qquad \textbf{(C)}\ 2\sqrt{6} \qquad \textbf{(D)}\ 5 \qquad \textbf{(E)}\ 2\sqrt{7}$

 \textbf{Answer: }C



\textbf{Solution 1}

Note as above that ABCD must be tangential to obtain the circle with maximal radius. Let $E$, $F$, $G$, and $H$ be the points on $AB$, $BC$, $CD$, and $DA$ respectively where the circle is tangent. Let $\theta=\angle BAD$ and $\alpha=\angle ADC$. Since the quadrilateral is cyclic(because we want to maximize the circle, so we set the quadrilateral to be cyclic), $\angle ABC=180^{\circ}-\alpha$ and $\angle BCD=180^{\circ}-\theta$. Let the circle have center $O$ and radius $r$. Note that $OHD$, $OGC$, $OFB$, and $OEA$ are right angles. 

 Hence $FOG=\theta$, $GOH=180^{\circ}-\alpha$, $EOH=180^{\circ}-\theta$, and $FOE=\alpha$. 

 Therefore, $AEOH\sim OFCG$ and $EBFO\sim HOGD$. 

 Let $x=CG$. Then $CF=x$, $BF=BE=9-x$, $GD=DH=7-x$, and $AH=AE=x+5$. Using $AEOH\sim OFCG$ and $EBFO\sim HOGD$ we have $r/(x+5)=x/r$, and $(9-x)/r=r/(7-x)$. By equating the value of $r^2$ from each, $x(x+5)=(7-x)(9-x)$. Solving we obtain $x=3$ so that $r=\boxed{\textbf{(C)}\ 2\sqrt{6}}$.

 



\textbf{Solution 2}

To maximize the radius of the circle, we also need to maximize its area. To maximize the area of the circle, the quadrilateral must be tangential (have an incircle). In a tangential quadrilateral, the sum of opposite sides is equal to the semiperimeter of the quadrilateral. $14+7=12+9$, so this particular quadrilateral has an incircle. By definition, given $4$ side lengths, a cyclic quadrilateral has the maximum area of any quadrilateral with those side lengths. Therefore, to maximize the area of the quadrilateral and thus the incircle, we assume that this quadrilateral is cyclic. 

 For cyclic quadrilaterals, \href{https://artofproblemsolving.com/wiki/index.php?title=Brahmagupta\%27s\_formula}{Brahmagupta's formula} gives the area as $\sqrt{(s-a)(s-b)(s-c)(s-d)}$ where $s$ is the semiperimeter and $a, b, c,$ and $d$ are the side lengths. Breaking it up into $4$ triangles, we see the area of a tangential quadrilateral is also equal to $r*s$. Equate these two equations.Substituting $s$, the semiperimeter, and $A$, the area and solving for $r$,we get $\boxed{\textbf{(C)}\ 2\sqrt{6}}$.

 



\textbf{Solution 3 (Trigonometry)}

By Pitot's Theorem, since $AB+CD=AD+BC$, there exists a circle tangent to all four sides of quadrilateral $ABCD$. Thus, we need to find the radius of this circle.

 Let the circle be tangent to $AB$, $BC$, $CD$, and $AD$ at points $P$, $Q$, $R$, and $S$, respectively. Also, let $AP=x$. Then $AS$ also equals $x$. Let the center of the circle be $O$. Observe that $AO$ bisects angle $\angle A$, so $\cot\frac{A}{2}=\frac{x}{r}$. Moreover, $\cot\frac{C}{2}=\frac{9-(14-x)}{r}=\frac{x-5}{r}$. But $\angle \frac{C}{2}=90^\circ-\angle \frac{A}{2}$, so $\cot\frac{C}{2}=\tan\frac{A}{2}$, and we find that $\frac{r}{x}=\frac{x-5}{r}$. Hence, $r^2=x(x-5)$.

 Similarly, $\cot\frac{B}{2}=\frac{14-x}{r}$, and $\cot\frac{D}{2}=\frac{12-x}{r}=\tan\frac{B}{2}=\frac{r}{14-x}$. Therefore, $r^2=\frac{12-x}{14-x}$. Putting this together with the above equation yields 
 \[(12-x)(14-x)=x(x-5)\]
 \[\Longrightarrow x^2-26x+168=x^2-5x\]
 \[\Longrightarrow 21x=168\]
 \[\Longrightarrow x=8.\] Thus, $r^2=8(8-5)=8(3)=24$, and $r=\boxed{\textbf{(C)}\ 2\sqrt{6}}$.

 


 



\textbf{Solution 4 (wowe)}

Let $r$ be the inradius. Then, the area of the quadrilateral can be expressed as $21r$. Clearly, the larger the area of the quadrilateral, the larger the value of the inradius.

 
Know that the area of a quadrilateral is maximized when it is cyclic. Applying Brahmagupta's, the area will be $\sqrt{(9)(7)(12)(14)} = 42\sqrt{6}$. Since $21r = 42\sqrt{6}$, $r = \boxed{\textbf{(C) } 2\sqrt6}$.

 -skibbysiggy

 



\textbf{Solution 5}

Why are the above solutions so yap buh.

 Assume that the quadrilateral is cyclic, because that is how the area fo a quadrilateral with fixed side lengths is maimised.

 Recall form googling random geometry theorems for AIME that the area of a bicentric quadrilateral with side lengths $a$, $b$, $c$, $d$, is $\sqrt{abcd}$, and this might seem too good to be true, but it is indeed true, and now we proceed with the solution.

 By Pitot's Theorem, since $AB+CD=AD+BC$, there exists a circle tangent to all four sides of quadrilateral $ABCD$. Thus, we need to find the radius of this circle. Note that $A = rs$ also holds for tangential quadrilaterals, and thus, the area of this tangential quarilateral which is known to be $42\sqrt{6}$, by $\sqrt{abcd}$ formula, has:

 
$rs = 42\sqrt{6}$
$r \cdot 21 = 42\sqrt{6}$
$r = \boxed{2\sqrt{6}}$, so $\boxed{(C)}$ $\square{}$

 -cheltstudent

 


	\item (\textit{AMC 12 12A 2011}) Triangle $ABC$ has $\angle BAC = 60^{\circ}$, $\angle CBA \leq 90^{\circ}$, $BC=1$, and $AC \geq AB$. Let $H$, $I$, and $O$ be the orthocenter, incenter, and circumcenter of $\triangle ABC$, respectively. Assume that the area of pentagon $BCOIH$ is the maximum possible. What is $\angle CBA$?


 $\textbf{(A)}\ 60^{\circ} \qquad \textbf{(B)}\ 72^{\circ} \qquad \textbf{(C)}\ 75^{\circ} \qquad \textbf{(D)}\ 80^{\circ} \qquad \textbf{(E)}\ 90^{\circ}$

 \textbf{Answer: }D



\textbf{Solution 1 (MAA)}

By the Inscribed Angle Theorem, 
 \[\angle BOC = 2\angle BAC = 120^\circ .\]Let $D$ and $E$ be the feet of the altitudes of $\triangle ABC$ from $B$ and $C$, respectively. In $\triangle ACE$ we get $\angle ACE = 30^\circ$, and as exterior angle 
 \[\angle BHC = 90^\circ + \angle ACE = 120^\circ .\]Because the lines $BI$ and $CI$ are bisector

 


\end{enumerate}